\documentclass[../../main.tex]{subfiles}
\begin{document}

% 年度の大きな表示
\begin{center}
  {\fontsize{48pt}{58pt}\selectfont \textbf{2006}\normalsize\textbf{年度}}
\end{center}

\vspace{10mm}

% 本文


\vspace{15mm}

% 出題テーマと難易度の表
\noindent\textbf{出題テーマと難易度}

\vspace{5mm}

\begin{center}
  \shadowbox{%
    \begin{tabular}{|c|p{8cm}|c|}
      \hline
      \textbf{問題番号} & \textbf{テーマ} & \textbf{難易度} \\
      \hline
      第1問 & 定積分の評価と不等式（$|\sin x|$の積分・ジョルダンの不等式） & \\
      \hline
      第2問 & 対数関数を含む不等式条件と領域の面積 & \\
      \hline
      第3問 & 平面図形（3円板の和集合の面積の最大化） & \\
      \hline
      第4問 & 空間ベクトル（四面体の中点と体積） & \\
      \hline
    \end{tabular}%
  }
\end{center}

\vspace{15mm}

% 解答
\noindent\textbf{解答}

\vspace{5mm}

\begin{center}
  \shadowbox{%
    \renewcommand{\arraystretch}{1.6}
    \begin{tabular}{|c|c|p{8cm}|}
      \hline
      \textbf{問題番号} & \textbf{小問} & \textbf{答え} \\
      \hline
      \multirow{3}{*}{第1問} & (1) & $I(n)=n$ \\
      \cline{2-3}
                             & (2) & 証明問題 \\
      \cline{2-3}
                             & (3) & 証明問題 \\
      \hline
      \multirow{3}{*}{第2問} & (1) & $t=\left(\dfrac{b}{a}\right)^{1/a}$で極小値$\dfrac{1}{a}-\dfrac{1}{a}\log\dfrac{b}{a}$ \\
      \cline{2-3}
                             & (2) & 図示（$D:0<x<y\le xe^{1-mx}$） \\
      \cline{2-3}
                             & (3) & $S=\left(e-\dfrac{5}{2}\right)\dfrac{1}{m^2}$ \\
      \hline
      第3問 & - & $\max S=2\pi+\dfrac{3\sqrt{3}}{2}$ \\
      \hline
      \multirow{3}{*}{第4問} & (1) & 証明問題 \\
      \cline{2-3}
                             & (2) & 証明問題 \\
      \cline{2-3}
                             & (3) & $\dfrac{\sqrt{2}}{12}$ \\
      \hline
    \end{tabular}%
    \renewcommand{\arraystretch}{1.0}
  }
\end{center}

\vspace{15mm}

% 解答時間
\noindent\textbf{試験形態}

\vspace{5mm}

\begin{center}
  \shadowbox{%
    \begin{tabular}{p{8cm}}
      （要確認）
    \end{tabular}%
  }
\end{center}

\end{document}
